\section{Elicitation 2026.10.15}

\begin{tabular}{ll}
\textbf{Date:} & 2026.10.15 \\
\textbf{Consultant:} & Hector Hill / Gy'ax \\
\textbf{Elicitors:} & Zahra, Andrew, Matthias \\
\textbf{Location:} & TFS elicitation room \\
\textbf{Context:} & Fourth elicitation for LING 431/531 Field Methods \\
\textbf{Audio:} & \texttt{October\_15\_2026\_HH.WAV} \\
\end{tabular}

\subsection{Andrew [00:00-24:42]}

[1:25]
\begin{exe}
\ex \gll k'sgoog̱ dim sidaam-'m, dim t'oyaxsi-'m simoogit sa lax̱ha\\
first \Fut{} start-1Pl.\Ii{}, \Fut{} thank-1Pl.\Ii{} creator day sky\\
'before we start, we thank our creator for the day'
\end{exe}

[2:11] --> asking about an itinerary on a trip
\begin{exe}
\ex \begin{xlist}
\item \gll * k'sgoog̱ dim Vancouver\\
{} first \Fut{} Vancouver \\
\item \gll Dim k'sgoog̱ fim Vancouver\\
\Fut{} first \Fut{} Vancouver\\
'First, Vancouver'
\end{xlist}
\end{exe}

[3:10]
\begin{exe}
\ex \gll Dim k'sgoog̱ dim Vancouver, dim-ii tl'om ga'a=hl Victoria, dim-ii hlisxw ga'a=hl Nanaimo\\
\Fut{} first \Fut{} Vancouver, \Fut{}-and ?? \Loc{}=\Cn{} Victoria, \Fut{}-and finished \Loc{}=\Cn{} Nanaimo\\
'First Vancouver, then Victoria, and finally Nanaimo'
\end{exe}

These next few are very tough because Henry was unable to help, Hector corrected the sentences into new forms on the 22nd, and I can't seem to find any translations in the literature...

[3:58] --> HH translates \textit{hagwin dulbin ju} as "it's close", but what the \textit{ju} actually is is anyone's guess
\begin{exe}
\ex \begin{xlist}
\item \gll hagwin dulbin ju=hl x̱biist ii sgi=hl dim=ist xwit sdo'oks ii g̱alen=hl g̱aldim aks\\
close approach ??=\Cn{} box and place.v=\Cn{} \Fut{}=\Qudd{} one.who.writes beside and after=\Cn{} container water\\
'The closest thing is the box, and after is the marker [== one who writes] and last is the water bottle'
\item hagwin dulbin ?? nit g̱o'o-'y, jag̱a-yex 'niin, jag̱a-yex 'nit\\
close approach ?? 3.\Iii{} \Loc{}-1.\Ii{},  across-?? 2.\Iii{}, across-?? 3.\Iii{}\\
'He's [== Matthias] the closest to me, you're across [the table], they're [== Zahra] across [the table]'
\item ii wa-'y tu'w-'n an \\
and find-1.\Ii{} far-2.\Ii{} ??\\
'And you're the furthest'
\end{xlist}
\end{exe}

[5:55]
\begin{exe}
\ex \begin{xlist}
\item \gll xstaa 'nisi-'m aa?\\
win 2Pl.\Iii{} \Q{}?\\
'Did you guys win?'
\item \gll Nee, moja xsta-'m\\
No, almost win-1Pl.\Ii{}\\
'No, we almost won'
\item \gll Nee-dii am=hl moja\\
\Neg{}-\Foc{} good=\Cn{} almost\\
'almost is not good'
\end{xlist}
\end{exe}

[6:30] --> I think this is in the context of a family (ie. "we're almost finished as a family"), but feel free to correct me
\begin{exe}
\ex \gll moja hlis-xw=s\\
almost finish-\Val{}=\Pn{}\\
'almost finished'
\end{exe}

[8:00] --> Nanaimo is closer to Victoria than Vancouver
\begin{exe}
\ex \begin{xlist}
\item \gll Moja dulbin ja Nanaimo e=hl Victoria ii nakw=hl Vancouver\\
almost short \Irr{} Nanaimo \Prep{}=\Cn{} Victoria and long=\Cn{} Vancouver\\
'Nanaimo is close to Victoria and long to Vancouver'
\item \gll nakw=hl Vancouver e=hl Victoria g̱an Nanaimo, ii moja dulbin ja=hl Nanaimo g̱an Victoria\\
long=\Cn{} Vancouver \Prep{}=\Cn{} Victoria \Pcnj{} Nanaimo, and almost short \Irr{}=\Cn{} Nanaimo \Pcnj{} Victoria\\
'Vancouver is long from Victoria and Nanaimo, and Nanaimo and Victoria are close' 
\end{xlist}
\end{exe}

[10:00]
\begin{exe}
\ex \begin{xlist}
\item \gll Moja nee-dii am sun loo'\\
almost \Neg{}-\Foc{} good this 1.\Prep{}\\
'This is not good for me'
\item \gll Moja nee-dii am-t\\
almost \Neg{}-\Foc{} good-3.\Ii{}\\
'It's not good'
\end{xlist}
\end{exe}

[11:00] --> Q: \textit{"How long will this trip be?"}
\begin{exe}
\ex \begin{xlist}
\item \gll Moja k'i hlox̱s\\
almost one month\\
'almost one month' - \textit{suggested answer}
\item \gll Moja k'i hlox̱s dim wil-'m\\
almost one month \Fut{} be-3Pl.\Ii{}\\
'It will be almost one month' - \textit{volunteered answer}
\end{xlist}
\end{exe}

[13:40] Q: "How many chairs will we need"
\begin{exe}
\ex \gll Dim=is g̱abi=hl ha'nii.ta'a=hl dim asag̱-an\\
\Fut{}=\Qudd{} amount=\Cn{} chair=\Cn{} \Fut{} want-2.\Ii{}\\
'How many chairs do you want?'
\end{exe}

HH: "\textit{dim asag̱an} is "something to add on", not "what you want""

AS: FV translates \textit{asag̱} as "want" pretty consistently... 

[14:50]
\textit{we had 5 people coming for dinner. I only had 2 chairs and you added 3 more}.
\begin{exe}
\ex \gll k'iinam-'y gilbil ha'nii ta'a loo'n  ii heksamox dim-'n heksamox gitnam gwile' ha'nii ta'a loo'n\\
give-1.\Ii{} two chair {} 2.\Prep{} \Ccnj{} again \Fut{}-1.\I{} again give\{1.\Ii{}\} three chair {} 2.\Prep{}\\
'I gave you 2 chairs and I'll give you three more' \textit{VG}
\end{exe}

[15:00]
\begin{exe}
\ex \begin{xlist}
\item \gll Geix gilbil ha'nii.ta'a asag̱an, dim-'n an gitnam=hl gwile'l loon\\
?? two chair want-2.\Ii{}, \Fut{}-1.\I{} ?? give=\Cn{} three 2.\Prep{}\\
'Here's the two chairs you wanted, but here's the three more that I'm gonna give you'
\item \gll Dim-'n heksamox k'ii'nam=hl gwile'l loo'n\\
\Fut{}-1.\I{} add give=\Cn{} three 2.\Prep{}\\
'I'll add three more'
\end{xlist}
\end{exe}

[16:25]
\begin{exe}
\ex \begin{xlist}
\item k'i g̱an=hl k'i ii sag̱ayt gilbil-t\\
 one \Pcnj{}=\Cn{} one and all.together two-3.\Ii{}\\
'One and one is two'
\item \gll k'i=hl g̱an=hl k'i ii sag̱ayt gilbil-t\\
one=\Cn{} \Pcnj{}=\Cn{} one and all.together two-3.\Ii{}\\
\textbf{Comment}: "This would mean there are 4 total"
\end{xlist}
\end{exe}

[19:23] --> Weird constructions, might be fully wrong, but the pieces all make sense kinda?
\begin{exe}
\ex \begin{xlist}
\item \gll gilbil=hl an aks-'m\\
two=\Cn{} \Ax{}? water-\Attr{}\\
\item \gll ii k'i=hl gii an aks ii isim\\
and one=\Cn{} ?? \Ax{} water and here\\
\end{xlist}
\end{exe}

[19:34]
\begin{exe}
\ex \gll ga'a sun\\
look.at this\\
'Look at this'
\end{exe}

The middle is extremely noisy and I can't make out what's being said

[19:40] HH: \textit{"When we're talking about your water bottles"}
\begin{exe}
\ex \gll dox̱=hl gilbil=hl g̱aldim aks isim xo'sos ii k'i=hl dii skit sdo'oks-'y\\
lie.on.PL=\Cn{} two=\Cn{} bottle water ?? ?? and one=\Cn{} \Foc{} here beside-1.\Ii{} \\
'Two water bottles are over there and there's one beside me'
\end{exe}

[23:30]
\begin{exe}
\ex \gll k'i g̱an=hl k'i an-jep=hl gilbil\\
one \Pcnj{}=\Cn{} one ??-make=\Cn{} two\\
'One and one makes two'
\end{exe}
No explicit comments here that there are 4 total, but this hasn't rigourously been tested.

[24:25]
\begin{exe}
\ex \gll na-di-k'i=hl\\
\Recp{}-\Com-one=\Cn{}\\
HH: \textit{"We wouldn't really use this in an addition context, we need to understand that there's two that makes it one"}
\end{exe}

[24:42]
\begin{exe}
\ex \gll Dim na-di-k'i 'nisi'm\\
\Fut{} \Recp{}-\Com{}-one 2Pl.\Iii{}\\
'Are you guys gonna be one [ie. married]'
\end{exe}

[25:25] --> Andrew's attempt at subtraction
\begin{exe}
\ex \gll gilbil jehl k'i jep=hl k'i\\
two lose.at one make=\Cn{} one\\
\textit{Intended: "Two losing one makes one"}
\end{exe}

[25:43]
\begin{exe}
\ex \gll dim saa guu dim k'i e=hl gilbil-t, dim-ix'm \tsb{k}'am k'i dim gina ski-t\\
\Fut{} away take \Fut{} one \Prep{}=\Cn{} two-3.\Ii{}, \Fut{}-?? only \Fut{} left lie.on-3.\Ii{}\\
'I'll take one of the two and there'll be only one left'
\end{exe}
This is following Rigsby p.332 on \textit{saa guu} as "away take"

\textit{\tsb{K}'am} is used in all leftover contexts, eg. 3-2 = "\tsb{k}'am gilbil"

[26:44] --> Not sure what the first word might be, possibly ??-\tsb{sexw} as in "go outside" based on his translation.
\begin{exe}
\ex \gll ajusexw sun gel\tsb{k} dim ii \tsb{k}'am gwile' loo'm\\
?? \Demprox{} outside \Fut{} and only three 1Pl.\Prep{}\\
'There's four of us and he goes outside, and there only three of us here'
\end{exe}

\subsection{Matthias [27:30-56:35]}

[28:00]: \textbf{Context:}  A surveyor comes to your door to make note of households that have 3 children or more for some government benefit. If the surveyor asks ``do you have three children?'', but you have five children, can you answer:

\begin{exe}
    \ex \glll \# gwlun g̱abiihl hlk'ii'y \\
    {} gwlun g̱abii=hl hlk'ii-'y \\
    {} three \Cnt{}=\Cn{} have.children-\Fsg{}.\Ii{} \\
    `I have three children (in total).'
\end{exe}

[28:45]: ``If you have five, you have to say the five'' e.g.:

\begin{exe}
    \ex \glll xwsdins g̱abii=hl hlk'ii-'y \\
    xwsdins g̱abii=hl hlk'ii-'y \\
    five \Cnt{}=\Cn{} have.children-\Fsg{}.\Ii{} \\
    `I have five children (in total).'
\end{exe}

[28:45]: What if you want to say that you have three children in the same context, to mean you have at least three children. In this case you can say the following. 

Inconclusive: idk if it was an audio error or Hector simply shook is head, I'm guessing the answer is no. But, this could be due to syntactic malformedness, so I probe further.

\begin{exe}
    \ex \glll * gwlunhl hlk'ii'y \\
    {} gwlun=hl hlk'ii-'y \\
    {} three=\Cn{} have.children-\Fsg{}.\Ii{} \\
    \textit{Intended: }`I have (at least) three children.'
\end{exe}

[29:55]:

\begin{exe}
    \ex \glll gwlunhl hlk'iinaa \\
    gwlun=hl hlk'ii-n=aa \\
    three=\Cn{} have.children-\Ssg{}.\Ii{}=\Q{} \\
    `Do you have three children?'
\end{exe}

[30:10]: \textbf{Context:} I have five kids, but I'm trying to say ``it's true I have three, but actually I have five''.

Less sure about the first gloss, but the second is repeated twice so it's pretty clear. Overall, it seems that \textit{g̱abii} leads to a ``total'' meaning, while the numeral in isolation is defeasible.

The \textit{iin k'ap} is necessary here (likely due to prgagmatic reasons).

\begin{exe}
    \ex
    \begin{xlist}
        \item \glll gwlunhl hlk'ii'y, iin k'ap xwsdins g̱abiihl hlk'ii  \\
    gwlun=hl hlk'ii-'y, ii=n k'ap xwsdins g̱abii=hl hlk'ii  \\
    three=\Cn{} have.children-\Fsg{}.\Ii{}, \Ccnj{}=\Fsg{}.\I{} \Verum{} five=\Cn{} \Cnt{}=\Cn{} have.children \\
    `I have three children, and actually I have five children (in total).'

    \item \glll ee, gwlunhl hlk'ii'y, iin k'ap xwsdinshl g̱abiidiit \\
    ee, gwlun=hl hlk'ii-'y, ii=n k'ap xwsdins=hl g̱abii-diit\\
    yes, three=\Cn{} have.children-\Fsg{}.\Ii{}, \Ccnj{}=\Fsg{}.\I{} \Verum{} five=\Cn{} \Cnt{}-\Tpl{}.\Ii{} \\
    `Yes, I have three childrem, and actually I have five.'
    \end{xlist}
\end{exe}

[32:18]: A surveyor comes to your door, and in this case you have to have three childeren to qualifiy for the policy. You do have three children. The following are possible responses.

The third one is ``almost like an elderly spokesperson'', but I'm guessing it is still grammatical. Unclear whether it is a full sentence however, since a plausible English reseponse could be ``Yes, three children''.

\begin{exe}
    \ex
    \begin{xlist}
    \item \glll gwlunhl g̱abiihl hlk'ii'y \\
    gwlun=hl g̱abii=hl hlk'ii-'y \\
    three=\Cn{} \Cnt{}=\Cn{} have.children-\Fsg{}.\Ii{} \\
    `I have three children (in total).' 
    \item \glll ee, gwilelhl g̱abiidiit \\
                ee, gwilel=hl g̱abii-diit \\
                yes, three=\Cn{} \Cnt{}-\Tpl{}.\Ii{} \\
            `Yes, there are three of them'
    \item \glll ee, gwlunhl hlk'ii'y \\
    ee, gwlun=hl hlk'ii-'y \\
    yes, three=\Cn{} have.children-\Fsg{}.\Ii{} \\
    `Yes, I have three children.'
    \end{xlist}
\end{exe}

[34:44]: This response is most appropriate since we are already talking about the children.

\begin{exe}
    \ex
    \begin{xlist}
    \item \glll gwlunhl hlk'iinaa \\
    gwlun=hl hlk'ii-n=aa \\
    three=\Cn{} have.children-\Ssg{}.\Ii{}=Q{} \\
    `Do you have three children?' 
    \item \glll nee, \tsb{k}'ap xwsdinshl g̱abiidiit \\
    nee, \tsb{k}'ap xwsdins=hl g̱abii-diit \\
    no, \Verum{} five=\Cn{} \Cnt{}.\Tpl{}.\Ii{} \\
    `No, there are actually five of them.'
    \end{xlist}
\end{exe}

[35:40]: This was brought up before, not sure what \textit{da'mis} is actually composed of.

\begin{exe}
    \ex \glll da'mis g̱abiidiit  \\
    da'mis g̱abii-diit  \\
    what(?) \Cnt{}-\Tpl{}.\Ii{} \\
    `How many of them (are there)?'
\end{exe}

[37:00]: \textbf{Context:} Zahra and I went on a walk. We then talk about the amount of bears that we saw. I saw two bears, Zahra saw three bears. 

I'm pretty sure there is a \textit{dii} \Foc{} in the third sentence (which Hector gave in answering this context), but not 100\% sure.

Also, I have no clue what the last word of the third gloss (starts at [37:50]). It sounds smething like ``\textit{glaat}'', but most likely it is \textit{ga'at} with third person series II agreement. Weird since it's supposedly embedded under \textit{ii}.

\begin{exe}
    \ex
    \begin{xlist}
    \item \glll gilbilhl smex ga'a'y \\
    gilbil=hl smex ga'a-'y\\
    two=\Cn{} bear see-\Fsg{}.\Ii{} \\
    `I saw two bears.' 
    \item \glll gwlelhl smex ga'an\\
    gwlel=hl smex ga'a-n\\
    three=\Cn{} bear see-\Ssg{}.\Ii{} \\
    `You saw three bears.' 
    \item \glll sag̱ayt xwsdinshl g̱abiihl smex ga'a'm, gilbilhl dii ga'a'y, ii gwlelhl ga'at\\
    sag̱ayt five=hl g̱abii=hl smex ga'a-'m, gilbil=hl dii ga'a-'y, ii gwlel=hl ga'a-t\\
    together five=\Cn{} \Cnt{}=\Cn{} bear see-\Fpl{}.\Ii{}, two=\Cn{} \Foc{} see-\Fsg{}.\Ii{}, \Ccnj{} three=\Cn{} see-\Tsg{}.\Ii{}  \\
    [Most likely]: `We saw five bears altogether, I saw two, and he saw three.' 
    \item \glll xwsdinshl g̱abiihl smex ga'a'm\\
    five=hl g̱abii=hl smex ga'a-'m\\
    five=\Cn{} \Cnt{}=\Cn{} bear see-\Fpl{}.\Ii{} \\
    `We saw (the same) five bears altogether.' 
    \end{xlist}
\end{exe}

[38:25]: If we don't have \textit{sag̱ayt}, we are ``not separating that you seen two and he seen three'', and in fact it means ``you both seen five bears at the same time, not separately''. With \textit{sag̱ayt}, it means ``you didn't see him see three bears, and he didn't see you see two bears''.

So, pretty clearly, we get a stronger distributive reading with \textit{sag̱ayt}, where we add separate bear-seeing occurences across two people, while wihtout it we are simply counting bear-seeing occurences across one (plural) individual. This definitely calls from some more thought.

We can also use \textit{sag̱ayt} if we saw the five bears together if we want to emphasize that there are two people. 

[42:00]: \textbf{Context:} Someone is chopping wood and you want to say that it's enough. 

Hector says \textit{hlaa} is `almost', but in a way this may reflect the fact that it means `it's good enough now', but it's used maybe politely to say `you have done enough now.' \textit{hlaa} reminds me of \textit{le} in Mandarin.

\begin{exe}
    \ex \glll hlaa hogax̱hl g̱abiihl jebin  \\
    hlaa hogax̱=hl g̱abii=hl jeb-n   \\
    \Incep{} good=\Cn{} \Cnt{}=\Cn{} do-\Ssg{}.\Ii{} \\
    `The amount you did is good.'
\end{exe}

[43:50]: \textbf{Context:} We want to bake a pie, and you are not sure if there are enough berries. 

Here the above construction doesn't really work, likey because \textit{maa'y} isn't a verb compared to \textit{jeb}. Instead Hector says:

\begin{exe}
    \ex
    \begin{exe}
        \item \glll amhl g̱abiihl maa'yaa  \\
    am=hl g̱abii=hl maa'y=aa   \\
    good=\Cn{} \Cnt{}=\Cn{} berry=\Q{} \\
    `Is the amount of berries good?'
    \item \glll amhl g̱abiihl maa'y  \\
    am=hl g̱abii=hl maa'y   \\
    good=\Cn{} \Cnt{}=\Cn{} berry \\
    `The amount of berries is good.'
    \end{exe}
\end{exe}

[46:25]: I'm really not sure what the full sentence here is, I'll have to revisit. Apparently it should mean `How much of this do you want?'.

\begin{exe}
    \ex \glll de g̱asg̱oot  \\
    de g̱asg̱oo-t  \\
    \Wh{}.\Spt{} \Mass{}-\Tsg{}.\Ii{}  \\
    `How much?'
\end{exe}

[46:42]: Rigsby dictionary p. 43 translates \textit{maadim} as snow (falling). I think it is used as a verb here however. Also, for what precedes it, it almost sounds like \textit{hats'im} `just' (e.g. Henry grammar p. 68), but the initial consonant sounds instead like a \textit{g̱}.

\begin{exe}
    \ex 
    \begin{xlist}
    \item \glll \tsb{k}'ay ts'im maadim  \\
    \tsb{k}'ay ts'im maadim \\
    just ?? snow  \\
    'It just snowed.' \textit{HH}
    \item \glll yukwhl maadim  \\
    yukw=hl maadim \\
    \Prog{}=\Cn{} snow  \\
    'It is snowing.'
     \item \glll hlaa dehl maadim  \\
    hlaa dehl maadim  \\
    \Incep{} ??? snow  \\
    Maybe: `It's starting to snow/snow is starting to come down.' 
    \end{xlist}
\end{exe}

[47:08]:

\begin{exe}
    \ex
    \begin{xlist}
    \item \glll * hilthl maaxws \\
    {} hilt=hl maaxws   \\
    {} deep=\Cn{} snow  \\
    Intended: `There is a lot of snow.'
    \item \glll hlephl maaxws \\
    hlep=hl maaxws   \\
    deep=\Cn{} snow  \\
    `There is deep snow.'
    \end{xlist}
\end{exe}

[48:15]:

\begin{exe}
    \ex \glll hlaa yukwhl maadim  \\
    hlaa yukw=hl maadim   \\
    \Incep{} \Prog{}=\Cn{} snow  \\
    `It's gonna snow.' \textit{HH}
\end{exe}

[48:48]: ``If you are bringing it in, there's buckets [...] if there's snow on your hat...'' then you can use the following. 

So this suggests that there is some sort of coercion possible, but as far as I can tell \textit{hilt} is only compatible with count nouns? 

\begin{exe}
    \ex \glll hilthl maaxws \\
    hilt=hl maaxws   \\
    many=\Cn{} snow  \\
    `There's a lot of snow.'
\end{exe}

[49:25]: TODO: Is this movement or is \textit{gilbilhl enjemhl} the predicate??

\begin{exe}
    \ex \glll gilbilhl enjemhl luu mitxwithl maaxws \\
    gilbil=hl enjem=hl luu mitxw-it=hl maaxws   \\
    two=\Cn{} pot=\Cn{} in full-\Sx{}=\Cn{} snow  \\
    `There's two pots that are filled with snow.' \textit{HH}
\end{exe}

[49:50]: The word is \textit{ts'ilim} but the first consonant sounds super voiced to me. I have simply no idea what the last word is - logic tells me it should be `put in', but I could not find any corresponding word anywhere.

\begin{exe}
    \ex \glll ts'ilim do\tsb{k}whl enjemhl maaxws inwin \\
    ts'ilim do\tsb{k}w=hl enjem=hl maaxws inwin    \\
    into bring=\Cn{} pot=\Cn{} snow ???   \\
    `Bring in the snow you put in the pots.' \textit{HH}
\end{exe}

[50:08]:

\begin{exe}
    \ex \glll hlaa yukwhl hlephl maaxws gelk \\
    hlaa yukw=hl hlep=hl maaxws gelk   \\
    \Incep{} \Prog{}=\Cn{} deep=\Cn{} snow outside \\
    `The snow is starting to pile up [=get deep?] outside' 
\end{exe}

[50:50]: \textbf{Context}: I was just outside and Andrew sees me and asks:

\begin{exe}
    \ex
    \begin{xlist}
        \item \glll * hindehl g̱asg̱oohl maaxws \\
    {} hinde=hl g̱asg̱oo=hl maaxws \   \\
    {} \Wh{}=\Cn{} \Mass{}=\Cn{} snow \\
    Intended: `How much snow is there?' 
        \item \glll hlephl maaxwsaa \\
    hlep=hl maaxws=aa  \\
    deep=\Cn{} snow=\Q{} \\
    `Is the snow deep?' 
    \item \glll \tsb{k}'ay yukwhl maadimaa \\
    \tsb{k}'ay yukw=hl maadim=aa  \\
    still \Prog{}=\Cn{} snow.fall=\Q{} \\
    `Is it still snowing?' 
    \end{xlist} 
\end{exe}

[52:00]: \textit{hlep 'yens} can be used for leaves, also:

\begin{exe}
    \ex \glll hlep dulpxw 'nii \\
    hlep dulpxw 'nii   \\
    deep short \Fsg{}.\Iii{} \\
    `I am short.' \textit{HH} 
\end{exe}

[52:49]: Hector explains that when someone is short it is like them being in the ground.

\begin{exe}
    \ex \glll hlep dulpxw gat tun \\
    hlep dulpxw get tun   \\
    deep short man \Pn{}=\Dem{}.\Prox{} \\
    `This man is short.' \textit{HH}
\end{exe}

[54:30]: TODO: I can't conclusively say what comes after \textit{naahl}. Will keep it unglossed for now.

But I'm pretty sure that in this case \textit{loo'm} is used, so it can almost be translated as ``Who is shorter \textit{to} us''.

\begin{exe}
    \ex \glll naahl g̱ay ts'a dulpxwt loo'm \\
    naa=hl g̱ay ts'a dulpxw-t loo-'m  \\
    who=\Cn{} ?? ?? short-\Tsg{}.\Ii{} \Obl{}-\Fpl{}.\Ii{} \\
    `Who is more shorter than us?' \textit{HH}
\end{exe}

[54:58]: I think when you want to say ``I am short'' in isolution you have to have \textit{hlep}. But wiht these contexts you don't have to have \textit{hlep}.

\begin{exe}
    \ex \glll g̱ay \tsb{k}'aa dulpxw 'niihl wegii'y \\
    g̱ay \tsb{k}'aa dulpxw 'nii=hl wegii-'y  \\
    \Cntr{} exceedingly short \Fsg{}.\Iii{}=\Cn{} brothers-\Fsg{}.\Ii{} \\
    `I am shorter than my brothers.' \textit{HH}
\end{exe}

[55:45]: \textit{nekwhl sint} is better for ``long summer''. 

\begin{exe}
    \ex \glll g̱ay yukwhl sint/maadim/gwooyim \\
    g̱ay yukw=hl sint/maadim/gwooyim  \\
    \Cntr{} \Prog{}=\Cn{} summer/winter/spring \\
    `Even now it's summer/winter/spring.' 
\end{exe}

\subsection{Zahra []}

[56:50]:
\begin{exe}
    \ex \glll
    k'i'yhl t’imḵ’aax g̱a’nekwihl ha’niitx̱ooḵxw \\
    k'i'yhl g̱a-’nakw=hl ha-’nii-tx̱ooḵ-xw \\
    five \Pl-foot=\Cn{} \Dist-long=\Cn{} \Ins-on-eat-\Val{} \\
    \glt `the table is one arm span (about six feet) long' (SF)
\end{exe}

[57:00]
\begin{exe}
    \ex \glll
    (*'wii) ’nekwhl ha’niitx̱ooḵxw \\
    (*'wii) ’nakw=hl ha’niitx̱ooḵxw \\
    big long=\Cn{} table \\
    \glt `long table' (SF)
\end{exe}

[discussion on above]

[57:30]
\begin{exe}
    \ex \glll
    nekw hlo' sild \\
    nekw hlo' sil-t \\
    long? go.\Pl? drunk-? \\
    \glt `you got drunk again?' (VF)
\end{exe}

HH: \textit{'nekw} has different meanings.

[58]:
\begin{exe}
    \ex \glll
\end{exe}